% Generated draft: requires mathematical and source review.
\documentclass[11pt]{article}
\usepackage[margin=25mm]{geometry}
\usepackage{fontspec}
\setmainfont{Latin Modern Roman}
\usepackage{amsmath,amssymb,mathtools,unicode-math}
\setmathfont{Latin Modern Math}
\usepackage{enumitem,needspace,xurl,hyperref}
\setlength{\parindent}{0pt}
\setlength{\parskip}{5pt}
\setlength{\emergencystretch}{3em}
\newcommand{\N}{\mathbb{N}}
\newcommand{\Z}{\mathbb{Z}}
\newcommand{\Q}{\mathbb{Q}}
\newcommand{\R}{\mathbb{R}}
\newcommand{\C}{\mathbb{C}}
\newcommand{\F}{\mathbb{F}}
\newcommand{\A}{\mathbb{A}}
\newcommand{\PP}{\mathbb{P}}
\newcommand{\E}{\mathbb{E}}
\newcommand{\eps}{\varepsilon}
\newcommand{\dd}{\,\mathrm{d}}
\begin{document}
\section*{UnsolvedMath ID: 10000002}
\textbf{Source problem code:} catalogue entry
\subsection{Definitions and mathematical statement}
Let \(W_n\) be the reduced words of length \(n\) over \(a,a^{-1},b,b^{-1}\); \(|W_n|=4\cdot3^{n-1}\). For a density \(d\), the source samples \(3^{nd}\) independent uniform relators from \(W_n\); an integer rounding convention is needed. Given the resulting sequence \(R_1,\ldots,R_m\), define the quotient chain
\[
G_k=\langle a,b\mid R_1,\ldots,R_k\rangle\qquad(0\le k\le m).
\]
On a realization containing a trivial quotient, put
\[
\tau_0=\min\{k:G_k=1\},\qquad
\tau_\infty=\max\{k<\tau_0:G_k\text{ is infinite}\}.
\]
The count requested is
\[
N_n=\left|\{[G_k]:\tau_\infty<k<\tau_0,\
1<|G_k|<\infty\}\right|,
\]
where brackets mean abstract group isomorphism class. Repeated unchanged quotients are counted once, not once per added relator.

The question asks for a useful upper bound on this finite-window count and conjectures an almost-sure bound uniform in \(n\). On a specified joint probability space, one possible interpretation is
\[
\Pr\!\left(\sup_nN_n<\infty\right)=1.
\]
A deterministic uniform constant, a random almost-sure bound, and tightness/high-probability bounds are distinct; the source does not choose between them or specify the coupling across \(n\).

There is an essential parity qualification. If \(n\) is even, sending both \(a,b\) to \(1\in\Z/2\Z\) kills every length-\(n\) relator, so every \(G_k\) still surjects onto \(\Z/2\Z\). Such a chain cannot collapse to the trivial group. Also a fixed finite sample of relators need not contain a collapse at its chosen density. Restricting to odd lengths, changing the terminal group, conditioning on collapse, or extending an infinite relator stream are different repairs; none is silently adopted. Until these choices are supplied, \(N_n\) is only conditionally defined by the displayed stopping indices.

The supplied summary reports no verified finite-window census bound and distinguishes related random-group asymptotic results from this sequential count. The title's property-T motivation is retained without adding a property-T hypothesis or a proof of the conjecture.
\subsection{Short English statement}
Bound the number of distinct finite nontrivial quotients in the sequential random-relator collapse window, ideally almost surely uniformly in word length. Even-length parity, the sample horizon, and the meaning of the joint almost-sure bound need clarification.
\subsection{Sources}
\begin{itemize}
\item Local UnsolvedMath registry, permanent ID \(10000002\), exact source code \texttt{catalogue entry}, ``Finite groups in the property-T collapse window.'' Proposer Itai Benjamini, 2007.
\item Benjamini, \emph{A 1'st order phase transition via property T}, May 2007, Question 1, one-page note; \url{https://www.wisdom.weizmann.ac.il/~itai/hplast.pdf}. Recorded access 2026-07-29; manual transcription checked against archived text.
\item Recorded preservation copy: \url{https://arquivo.pt/wayback/20201231041601id_/https://www.wisdom.weizmann.ac.il/~itai/hplast.pdf}.
\item Supplied OPEN-TRIAGE summary mentions Ollivier and Kotowski--Kotowski as related density-model literature, but provides no exact theorem or primary citation resolving the finite-window count.
\end{itemize}
\end{document}
